\begin{blocksection}
    \question What would Python display? Draw box-and-pointer diagrams for the following:
    
    \begin{lstlisting}
    >>> a = [1, 2, 3]
    >>> a
    \end{lstlisting}
    \begin{solution}[.25in]
    \begin{lstlisting}
    [1, 2, 3]
    \end{lstlisting}
    \end{solution}
    
    \begin{lstlisting}
    >>> a[2]
    \end{lstlisting}
    \begin{solution}[.25in]
    \begin{lstlisting}
    3
    \end{lstlisting}
    \end{solution}
    
    \begin{lstlisting}
    >>> a[:2]
    \end{lstlisting}
    \begin{solution}[.25in]
    \begin{lstlisting}
    [1, 2]
    \end{lstlisting}
    \end{solution}

    \begin{lstlisting}
    >>> len(a)
    \end{lstlisting}
    \begin{solution}[.25in]
    \begin{lstlisting}
    3
    \end{lstlisting}
    \end{solution}

    \begin{lstlisting}
    >>> c = [x * 2 for x in a]
    >>> c
    \end{lstlisting}
    \begin{solution}[.25in]
    \begin{lstlisting}
    [2, 4, 6]
    \end{lstlisting}
    \end{solution}

    \begin{lstlisting}
    >>> sum(c)
    >>> max(c)
    >>> all(c)
    >>> all([0, 1, 2, 3])
    \end{lstlisting}
    \begin{solution}[.25in]
    \begin{lstlisting}
    12
    6
    True
    False
    \end{lstlisting}
    \end{solution}

    \begin{lstlisting}
    >>> b = [1, 2, 3, a, 4]
    >>> b
    \end{lstlisting}
    \begin{solution}[.25in]
    \begin{lstlisting}
    [1, 2, 3, [1, 2, 3], 4]
    \end{lstlisting}
    \end{solution}

    \begin{lstlisting}
    >>> b[3][2]
    \end{lstlisting}
    \begin{solution}[.25in]
    \begin{lstlisting}
    3
    \end{lstlisting}
    \end{solution}    
\end{blocksection}

    \begin{questionmeta}
        \textbf{Teaching Tips}
        \begin{enumerate}
                \item The goal of this problem is to introduce lists, box and pointer diagrams, list indexing and slicing, list comprehensions, aggregation functions, and nested lists.
                \item Encourage students to draw box and pointer diagrams if they seem stuck, but it is not necessary to draw a box and pointer diagram for all questions. 
        \end{enumerate}
    \end{questionmeta}

